% =========================================================
% Proof: Archimedean Property
% Source: volume-iii/analysis/bounding/notes/notes-completeness.tex
% Mode: standard
% =========================================================

\subsection*{Archimedean Property}
\label{prf:archimedean-property}

\begin{remark*}[Return]
$\leftarrow$ \hyperref[thm:archimedean-property]{Back to Theorem in Notes}
\end{remark*}

\begin{proof}
\Claim $\forall x \in \mathbb{R},\; \exists n \in \mathbb{N},\; n > x$.

\Strategy Contradiction.
Assume $\mathbb{N}$ is bounded above in $\mathbb{R}$.
The completeness axiom then gives $\sup \mathbb{N} \in \mathbb{R}$.
We exhibit an element of $\mathbb{N}$ exceeding this supremum,
contradicting it being an upper bound.

\medskip
Assume for contradiction that $\mathbb{N}$ is bounded above in $\mathbb{R}$.
Then in particular there exists $x \in \mathbb{R}$ that is an upper bound
for $\mathbb{N}$.
By the Axiom of Completeness, $s := \sup \mathbb{N}$ exists in $\mathbb{R}$.

By the $\varepsilon$-characterisation of the supremum, for every
$\varepsilon > 0$ there exists $n \in \mathbb{N}$ such that
\[
  n > s - \varepsilon.
\]
Setting $\varepsilon := 1$, there exists $n \in \mathbb{N}$ satisfying
\begin{equation}
\label{eq:arch-near-sup}
  n > s - 1.
\end{equation}
Adding $1$ to both sides of \eqref{eq:arch-near-sup}:
\[
  n + 1 > s.
\]
By the successor axiom, $n + 1 \in \mathbb{N}$.
But $n + 1 > s = \sup \mathbb{N}$ contradicts $s$ being an upper bound
for $\mathbb{N}$.

\Hence our assumption was false.
Therefore $\mathbb{N}$ is unbounded in $\mathbb{R}$:
for any $x \in \mathbb{R}$, $x$ is not an upper bound for $\mathbb{N}$,
so there exists $n \in \mathbb{N}$ with $n > x$. \AsReq
\end{proof}

\begin{remark*}[Proof shape]
Contradiction via the $\varepsilon$-characterisation of supremum.
If $\mathbb{N}$ were bounded above, completeness would give $s = \sup\mathbb{N}$.
The $\varepsilon$-characterisation with $\varepsilon = 1$ produces $n \in \mathbb{N}$
with $n > s - 1$, so $n + 1 > s$ — but $n+1 \in \mathbb{N}$ by the successor axiom,
contradicting $s$ being an upper bound.

\FlashProof{Contradiction. Assume $\mathbb{N}$ bounded above; completeness gives $s=\sup\mathbb{N}$. $\varepsilon$-characterisation with $\varepsilon=1$ gives $n\in\mathbb{N}$ with $n>s-1$, so $n+1>s$. By successor axiom $n+1\in\mathbb{N}$, contradicting $s$ being an upper bound.}
\FlashUses{Axiom of Completeness; $\varepsilon$-characterisation of supremum; successor axiom.}
\end{remark*}

\begin{remark*}[Why completeness is essential]
The Archimedean property does \emph{not} follow from the ordered field axioms alone.
Non-Archimedean ordered fields exist (e.g.\ the field of formal Laurent series).
Completeness is what rules them out for $\mathbb{R}$ --- without it,
$\mathbb{N}$ could be bounded above with no contradiction.
\end{remark*}

\begin{remark*}[Dependencies]
\begin{itemize}
  \item \hyperref[axiom:completeness]{Axiom of Completeness}
  \item \hyperref[thm:supremum-approximation]{$\varepsilon$-characterisation of supremum}
  \item Successor axiom: $n \in \mathbb{N} \Rightarrow n+1 \in \mathbb{N}$
\end{itemize}
\end{remark*}
